\documentclass{article}
\usepackage{graphicx} % Required for inserting images
\usepackage[margin=1in]{geometry}
\usepackage{amsmath}
\usepackage{amsfonts}
\usepackage{amsthm}
\usepackage{thm-restate}
\usepackage[]{hyperref}
\usepackage{cleveref}
\usepackage{xcolor}
\usepackage{bbm}
\usepackage{dirtytalk}
\usepackage{graphicx}
\usepackage{subcaption}
\usepackage{float}
\usepackage{tikz}
\usepackage{algorithm}
\usepackage[noend]{algpseudocode}
\usetikzlibrary{arrows.meta,decorations.markings, calc, positioning}


\declaretheorem[name=Theorem,numberwithin=section]{thm}
\declaretheorem[name=Lemma,numberwithin=section]{lemma}
\declaretheorem[name=Assumption,numberwithin=section]{asmpt}
\declaretheorem[name=Proposition,numberwithin=section]{prop}
\declaretheorem[name=Corollary,numberwithin=section]{corollary}
\declaretheorem[name=Remark,numberwithin=section]{rmk}
\declaretheorem[name=Definition,numberwithin=section]{defn}

\newcommand{\man}{\mathcal{M}}

\definecolor{cyan}{cmyk}{1, 0.4, 0, 0} 
\newcommand{\ndfg}[1]{{\color{magenta}[Florian: #1]}}
\newcommand{\tls}[1]{{\color{cyan}[TLS: #1]}}
\newcommand{\pac}{\mathcal{P}_{2, \text{ac}}(\mathbb{R}^d)}
\newcommand{\posm}{\mathfrak{M}_+(\mathbb{R}^d)}

\usepackage[numbers]{natbib}
\bibliographystyle{unsrtnat}

\title{Hellinger Kantorovich Parallel Transport for Counterfactual Dynamics Estimation}
\author{Florian Gunsilius, Gonzalo Mena, Tristan Saidi, Larry Wasserman}

\begin{document}

\maketitle


\section{Introduction}

\section{Wasserstein and HK Spaces}

In this section we describe the construction of Wasserstein space and its pseudo-Riemannian properties.  We will then use this construction to define an analogous space corresponding to general positive measures instead of probability measures based on the principles of unbalanced transport. This construction follows closely the construction described in \cite{clancy2021interpolating, chewi2024statisticaloptimaltransport, liero2018optimal, liero2016optimal, chizat2018interpolating, kondratyev2016newoptimaltransportdistance}.  

\subsection{Optimal Transport}
Define $\mathcal{P}_{2}(\mathbb{R}^d)$ to be the set of probability measures over $\mathbb{R}^d$ with finite $2^{\text{nd}}$ moment. The $2$-Wasserstein distance between two probability measures $\mu, \nu \in \mathcal{P}_{2}(\mathbb{R}^d)$ is defined by
\begin{equation}
    W_2(\mu, \nu) = \inf_{\gamma \in \Gamma_{\mu, \nu}}\left(\int\|x - y\|_2^2\gamma(dx, dy)\right)^{1/2} \label{eq: 2 wasserstein dist}
\end{equation}
where $\Gamma_{\mu, \nu}$ denotes the set of couplings of $\mu$ and $\nu$. The $2$-Wasserstein distance (which we will henceforth refer to as the Wasserstein distance) is indeed a metric, which renders $(\mathcal{P}_{2, \text{ac}}(\mathbb{R}^d), W_2)$ a metric space. A key result is that of Brenier, who showed that the optimal coupling takes the form $(X, \nabla \varphi(X)), X \sim \mu$ for some convex $\varphi$ when $\mu$ is absolutely continuous with respect to the Lebesgue measure.

\begin{restatable}[Brenier]{thm}{}
    Let $\mu, \nu \in \mathcal{P}_2(\mathbb{R}^d)$ be probability measures such that $\mu$ has a density, and let $X \sim \mu$. If $\gamma^*$ is optimal for \cref{eq: 2 wasserstein dist}, then there exists a convex function $\varphi:\mathbb{R}^d \rightarrow \mathbb{R}$ such that $(X, \nabla\varphi(X))\sim \gamma^*$. 
\end{restatable}

This theorem guarantees that if the source measure has a density, then the optimal transport coupling can be written as the source measure $\mu$ and $\nabla \varphi_\#\mu$, the pushforward of the source under some convex function. The function $\nabla\varphi$ is often referred to as the \textit{Brenier map}. This key theorem will be central to our construction of pseudo-Riemannian structure on the space of probability measures. 

\subsection{Dynamic Formulation and Unbalanced Transport}

Now we'll define $\mathcal{P}_{2, \text{ac}}(\mathbb{R}^d)$ to be the set of absolutely continuous probability measures over $\mathbb{R}^d$ with finite second moment. Let $(v_t)_{t \geq 0}$ be a time dependent family of vector fields over $\mathbb{R}^d$, and consider the ODE $\dot X_t = v_t(X_t)$. Let $\mu_t$ denote the law of $X_t$, where $X_0 \sim \mu_0$ and $X_t$ evolves according to the ODE described previously. Then, the dynamics of $\mu_t$ obey the so-called \textit{continuity equation},
\begin{equation}
    \partial_t\mu_t + \nabla \cdot (\mu_t v_t) = 0 \label{eq: balanced continuity eqn}
\end{equation}
in the weak sense. Namely, for all compactly supported and smooth test functions $\varphi: \mathbb{R}^d \rightarrow \mathbb{R}$, we have
\[\partial_t \int \varphi d\mu_t = \int \langle \nabla \varphi, v_t\rangle d\mu_t.\]
It turns out, one can identify optimal transport maps with vector fields satisfying the continuity equation that are optimal in a certain sense. This gives rise to the \textit{dynamic} formulation of optimal transport.

\begin{restatable}[Benamou-Brenier]{thm}{}
    Let $\mu_0, \mu_1 \in \mathcal{P}_{2, \text{ac}}(\mathbb{R}^d)$. Then 
    \begin{equation*}
        W_2(\mu_0, \mu_1) = \inf \left\{\left(\int_{0}^1\|v_t\|^2_{L_2(\mu_t)} dt\right)^{1/2} \, \Bigg| \, (\mu_t, v_t) \text{ solves \cref{eq: balanced continuity eqn}}\right\}.
    \end{equation*}
    Moreover, the optimal curve $(\mu_t)_{t\geq 0}$ is unique and is described by $X_t \sim \mu_t$, where $X_t = (1-t)X_0 + tX_1$ and $(X_0, X_1) \sim \gamma^* \in \Gamma_{\mu_0, \mu_1}$ with $\gamma^*$ being an optimal coupling.
\end{restatable}

This dynamical formulation, which considers minimum cost transport between probability measures, can be augmented to allow for \textit{unbalanced} transport -- transport where the source and target measures do not necessarily have total mass $1$. To do this, we can create a modified continuity equation that allows for the creation and destruction of mass,
\begin{equation}
    \partial_t\mu_t + \nabla \cdot (\mu_tv_t) = 4\alpha_t\mu_t. \label{eq: unbalanced continuity eqn}
\end{equation}
This leads to an analogous notion of distance between general non-negative measures, commonly referred to as the \textit{Wasserstein-Fisher-Rao} (WFR) or the \textit{Hellinger-Kantorovich} (HK) distance. In fact, in \citet{liero2018optimal} it was shown that HK is indeed a metric on the space of non-negative measures over $\mathbb{R}^d$. 

\begin{restatable}[Hellinger Kantorovich Distance]{defn}{}
    Let $\mathfrak{M}_+(\mathbb{R}^d)$ denote the space of non-negative measures on $\mathbb{R}^d$. For any $\mu_0, \mu_1 \in \mathfrak{M}_+(\mathbb{R}^d)$ we define the Hellinger-Kantorovich distance to be
    \begin{equation}
        \text{HK}(\mu_0, \mu_1) = \inf\left\{\int_0^1 \left(\|v_t\|_{L_2(\mu_t)}^2 + |\alpha_t|_{L_2(\mu_t)}^2\right)dt \, \Bigg| \, (\mu_t, v_t,  \alpha_t) \text{ solves \cref{eq: unbalanced continuity eqn}}\right\}.
    \end{equation}
\end{restatable}

\subsection{Riemannian structure of $(\mathcal{P}_{2, \text{ac}}(\mathbb{R}^d), W_2)$ and $(\mathfrak{M}_+(\mathbb{R}^d), \text{HK})$}

\textbf{The Tangent Space}. We are almost in a position to describe the Riemannian-like structure of $(\mathcal{P}_{2, \text{ac}}(\mathbb{R}^d), W_2)$ and $(\mathfrak{M}_+(\mathbb{R}^d), \text{HK})$. But first, we will define a notion of a \say{derivative} in a general metric space.

\begin{restatable}[Metric Derivative]{defn}{}
    Let $(\mathcal{X}, d)$ be a metric space and $(x_t)_{t \geq 0}$ be a curve in $\mathcal{X}$. The metric derivative of the curve at time $t$ is given by
    \begin{equation*}
        |\dot x|(t) := \lim_{s\rightarrow t}\frac{d(x_s, x_t)}{|s - t|}
    \end{equation*}
    provided that the limit exists. 
\end{restatable}

Now, for a pair of probability measures $\mu, \nu \in \mathcal{P}_{2, \text{ac}}(\mathbb{R}^d)$ where $\mu$ admits a density, we will write $T_{\mu \rightarrow \nu}$ for the Brenier map from $\mu$ to $\nu$ and we will write $|\dot \mu|$ for the metric derivative of a curve in $\mathcal{P}_{2, \text{ac}}(\mathbb{R}^d)$ with respect to the Wasserstein metric.

\begin{restatable}[\citet{chewi2024statisticaloptimaltransport}]{thm}{}
    Let $(\mu_t)_{t\geq 0}$ be a regular curve, i.e. $\mu_t$ admits a density and $|\dot \mu_t|$ exists for all $t \geq 0$. Then for every family of vector fields $(v_t)_{t\geq 0}$ for which \cref{eq: balanced continuity eqn} holds, it holds that $|\dot \mu_t|(t) \leq \|v_t\|_{L_2(\mu_t)}$ for all $t \geq 0$. Moreover, there exists a unique family $(v_t)_{t\geq 0}$ such that \cref{eq: balanced continuity eqn} holds and for which $|\dot \mu_t|(t) = \|v_t\|_{L_2(\mu_t)}$ for every $t \geq 0$. This family is such that
    \begin{equation*}
        v_t = \lim_{h \rightarrow0^+}\frac{T_{\mu_t\rightarrow \mu_{t+h}}- I_d}{h} \qquad \text{ in $L_2(\mu_t)$}
    \end{equation*}
    where $I_d$ is the identity map. It also holds that 
    \begin{equation*}
        v_t = \min \left\{\|\tilde{v}_t\|_{L_2(\mu_t)}^2 \, \Big|\, \partial_t\mu_t + \nabla \cdot (\tilde v_t\mu_t) = 0\right\}.
    \end{equation*}
\end{restatable}

Note that Brenier's theorem indicates that the velocity field that coincides (in an $L_2(\mu_t)$ sense) with the metric derivative is a limit of gradients -- in particular, $T_{\mu_t \rightarrow \mu_{t+h}} - I_d = \nabla (\varphi - \frac{1}{2}\|x\|_2^2)$. This key fact gives rise to an intuitive definition of the tangent space and a metric tensor for $(\mathcal{P}_{2, \text{ac}}(\mathbb{R}^d), W_2)$ at a measure $\mu$. 

\begin{restatable}[Wasserstein Tangent Space]{defn}{}
    Let $\mu \in \pac$. We define the tangent space to $\pac$ at $\mu$ to be 
    \begin{equation}
        T_\mu \pac = \overline{\left\{\nabla \psi \, | \, \psi \in \mathcal{C}_c^{\infty}\right\}}^{L_2(\mu)}
    \end{equation}
    where $\overline{\{\cdot \}}^{L_2(\mu)}$ denotes the $L_2(\mu)$ closure and $\mathcal{C}_c^{\infty}$ denotes the set of compactly supported and smooth maps from $\mathbb{R}^d \rightarrow \mathbb{R}$. We also endow $T_\mu \pac$ with the metric tensor
    \begin{equation}
        \langle \nabla \psi_1, \nabla \psi_2\rangle_{\mu} = \int \langle\nabla\psi_1, \nabla\psi_2 \rangle d\mu. 
    \end{equation}
\end{restatable}
We are now in a position to define analogous structure on the Hellinger Kantorovich space, $(\posm, \text{HK})$. Starting from the augmented continuity equation, one can show that the minimum norm solution is unique and is the gradient of a potential. We can then use this to construct the tangent space and the metric tensor as we did for the Wasserstein space.
\begin{restatable}[\citet{clancy2021interpolating}]{prop}{}
    We have
    \begin{equation*}
        (v_t, \alpha_t) = \min\left\{\|\tilde v_t\|_{L_2(\mu_t)^2} + |\tilde \alpha_t|_{L_2(\mu_t)^2  }  \, \Big| \, \partial_t\mu_t + \nabla \cdot (\tilde v_t \mu_t) = 4\tilde \alpha_t \mu_t\right\}
    \end{equation*}
    with $(v_t, \alpha_t) \in \overline{\left\{(\nabla \varphi, \varphi) \, | \, \varphi \in \mathcal{C}^{\infty}_{c}\right\}}^{L_2(\mu_t)}$.
\end{restatable}
\begin{restatable}[Hellinger-Kantorovich Tangent Space]{defn}{}
    Let $\mu \in \posm$. We define the tangent space to $\posm$ at $\mu$ to be 
    \begin{equation}
        T_\mu \posm = \overline{\left\{(\nabla \varphi, \varphi) \, | \, \varphi \in \mathcal{C}_c^{\infty}\right\}}^{L_2(\mu)}.
    \end{equation}
    We also endow $T_\mu \posm$ with the metric tensor
    \begin{equation}
        \langle s_1,  s_2\rangle_{\mu} = \int \left(\langle\nabla\varphi_1, \nabla\varphi_2\rangle + 4\varphi_1\varphi_2 \right) d\mu
    \end{equation}
    for any $s_1, s_2\in T_\mu \posm$, with $s_i = (\nabla \varphi_i, \varphi_i)$ and $\varphi_i \in \mathcal{C}_c^\infty$. 
\end{restatable}

\noindent \textbf{The Covariant Derivative.} 
On an abstract manifold $(\man, g)$ that isn't embedded in an ambient space, we have no obvious way to compare vectors in the tangent space at two points $p, q \in \man$, $p \neq q$. Therefore, we need a way of connecting the separate vector spaces $T_p\man$ and $T_q\man$ when $p \neq q$. One can achieve this by defining a rule $\nabla$ for differentiating vector fields against each other on $\man$ in a way that preserves the structure of the metric $g$.

\begin{restatable}[The Covariant Derivative, \citet{petersen2006riemannian}]{defn}{}
    For a manifold $(\man, g)$ we define the covariant derivative $\nabla$ to be a rule that assigns to each pair of vector fields $X$, $Y$ over $\man$ another vector field $\nabla_XY$ satisfying 
    \begin{enumerate}
        \item \textit{Linearity:} $\nabla_{aX+bZ}Y = a\nabla_XY + b\nabla_ZY.$
        \item \textit{Leibniz/derivation property:} for $f \in C^{\infty}(\man)$ we have $\nabla_X(fY) = X(f)Y + f\nabla_XY$, where $X(f) = g(X, \nabla f)$ is the derivative of $f$ in the direction $X$.
    \end{enumerate}
\end{restatable}

While this gives us a way to connect tangent spaces, the covariant derivative might distort the geometry induced by the metric $g$ -- i.e. it might not be \say{metric compatible}. A fundamental result of Riemannian geometry, however, is that for any Riemannian manifold $(\man, g)$ there exists a \textit{unique} connection that is torsion free and is \say{metric compatible}. 

\begin{restatable}[The fundamental theorem of Riemannian Geometry, \cite{petersen2006riemannian}]{thm}{}
    If $\man$ is endowed with a Riemannian metric $g$, then there exists a unique connection $\nabla$ called the \textit{Levi-Civita connection} that is
    \begin{enumerate}
        \item \textit{Torsion free:} $\nabla_XY - \nabla_YX = [X,Y]$, where $[X,Y]$ is the Lie bracket of $X$ and $Y$. 
        \item \textit{Metric compatible:} $Xg(Y, Z) = g(\nabla_XY, Z) + g(Y, \nabla_XZ).$ 
    \end{enumerate}
\end{restatable}
In the case of $(\pac, W_2)$ and $(\posm, \text{HK})$ we can explicitly construct the covariant derivative, but we need to manually verify that they are torsion free and metric compatible as neither spaces are true Riemannian Manifolds.

\begin{restatable}[Wasserstein Covariant Derivative, \citet{clancy2021interpolating}]{prop}{} \label{Wasserstein covariant deriv}
    Let $(\mu_t)_{t\geq 0}$ be a curve through $\pac$ with tangent field $\nabla \varphi_t$ solving \cref{eq: balanced continuity eqn}, and therefore driving the dynamics of $\mu_t$. Also let $v_t$ be another vector field along $\mu_t$ and let $\Pi_{\mu_t}$ be the orthogonal projection onto $T_{\mu_t}\pac$ in $L_2(\mu_t)$. Then the differential operator $\nabla_{\nabla\varphi_t}^{W_2}$ given by
    \begin{equation*}
        \nabla_{(\nabla\varphi_t)}^{W_2}v_t = \Pi_{\mu_t}\left(\partial_t v_t + \nabla v_t \cdot \nabla \varphi_t\right)
    \end{equation*}
    is a valid covariant derivative, is torsion free, and is metric compatible. 
\end{restatable}

\begin{restatable}[Hellinger-Kantorovich Covariant Derivative, \citet{clancy2021interpolating}]{prop}{} \label{HK covariant deriv}
    Let $(\mu_t)_{t\geq 0}$ be a curve through $\posm$ with tangent field $(\nabla \varphi_t, \varphi_t)$ solving \cref{eq: unbalanced continuity eqn}, and therefore driving the dynamics of $\mu_t$. Also let $(v_t, \beta_t)$ be another vector field along $\mu_t$ and let $\Pi_{\mu_t}$ be the orthogonal projection onto $T_{\mu_t}\posm$ in $L_2(\mu_t)$. Then the differential operator $\nabla_{(\nabla\varphi_t, \varphi_t)}^{\text{HK}}$ given by
    \begin{equation*}
        \nabla_{(\nabla\varphi_t, \varphi_t)}^{\text{HK}} \begin{pmatrix}
            v_t \\
            \beta_t
        \end{pmatrix} = \Pi_{\mu_t}\begin{pmatrix}
            \partial_t v_t + \nabla v_t \cdot \nabla \varphi_t + 2\beta_t\nabla\varphi_t \\
            \partial_t \beta_t + \frac{1}{2}\langle \nabla \beta_t, \nabla \varphi_t\rangle + 2\varphi_t\beta_t
        \end{pmatrix}
    \end{equation*}
    is a valid covariant derivative, is torsion free, and is metric compatible. 
\end{restatable}

\subsection{Hellinger-Kantorovich via Cone Lifting}

Amazingly, \citet{liero2016optimal} show that Hellinger-Kantorovich geometry arises from $W_2$ geometry on an augmented base domain. In particular, this construction entails lifting positive measures on the base space $\mathbb{R}^d$ to probability measures on $\mathfrak{C}(\mathbb{R}^d)$, the metric cone of the base space ($\mathbb{R}^d$). Under a particular set of lifts, it turns out that the Wasserstein geometry on the cone is identical to the Hellinger-Kantorovich geometry on the original space.

\begin{restatable}[Metric cone, \citet{burago2001course}]{defn}{}
    For a Riemannian manifold $(\man, g)$, the metric cone is
    \[\mathfrak{C}(\man) = \man \times \mathbb{R}_{+} / \{(x, 0) \sim (y, 0)\}\]
    where a point in $\mathfrak{C}(\man)$ is written as $(x, r)$. The metric tensor on the cone is given by $g_{(x,r)} = dr^2 + r^2g_x.$
\end{restatable} 
\tls{WIP}The metric (as one would have in a metric space) is given by
\begin{equation}
    d_{\mathfrak{C}}((x_0, r_0), (x_1, r_1))^2 = r_0^2 + r_1^2 - 2r_0r_1\cos(d_{\man}(x_0, x_1) \land \pi).
\end{equation}
Now, given a measure $\lambda \in M(\mathfrak{C})$ we can project it to a measure $\mathfrak{B}\lambda \in M(\mathbb{R}^d)$ with
\begin{equation}
    \int f(x)d\mathfrak{B}\lambda(x) = \int r^2 f(x) d\lambda(x, r) \qquad \text{for all test functions }f.
\end{equation}
Clearly, there are many possible ways to define a lift of a measure -- any $\lambda$ of the form
\begin{equation}
    \lambda(dx, dr) = \eta(x, dr)\mathfrak{B}\lambda(dx) \qquad \text{with } \qquad \int_0^\infty r^2\eta(x, dr) = 1 \text{ \quad for $\mathfrak{B}\lambda$ a.e. $x$}
\end{equation}
will do. To see this plug into the projection formula above and we obtain
\[\int r^2 f(x)\eta(x, dr)\mathfrak{B}\lambda(dx) = \int f(x) \left(\int r^2 \eta(x, dr) \right) \mathfrak{B}\lambda(dx) = \int f(x) \mathfrak{B}\lambda(dx) \]
as desired. \citet{liero2016optimal} show that one can characterize the Wasserstein-Fisher-Rao distance on a base space to the Wasserstein on the lifted cone space,
\begin{equation} \label{eq: cone wasserstein and hk equivalence}
    \text{HK}(\mu_0, \mu_1) = \min\left\{W_{\mathfrak{C}}(\lambda_0, \lambda_1) \,\Big|\, \mathfrak{B}\lambda_0 = \mu_0, \mathfrak{B}\lambda_1 = \mu_1\right\}.
\end{equation}

\section{Parallel Transport}

\textbf{Exact Parallel Transport via the Covariant Derivative.} Now we are in a position where we can define parallel transport. Intuitively, a vector field along a curve is \say{unchanging} loosely-speaking if its derivative is zero. Thus, in the context of abstract manifolds, we stipulate that a vector field along a curve is the parallel transport of a source vector if its covariant derivative along the curve is zero.
\begin{restatable}[\citet{lee2018introduction}]{defn}{} \label{parallel transport on general manifolds}
    Let $\man$ be a smooth Riemannian manifold. A smooth vector field $X$ along a smooth curve $\gamma$ is said to be parallel along $\gamma$ with respect to the Levi-Civita connection if $\nabla_{\dot \gamma}X \equiv 0$.
\end{restatable}
This characterization now allows us to define parallel transport on $(\pac, W_2)$ and $(\posm, \text{HK})$ using the covariant derivative. Consider a smooth curve $\mu_t$ through $\pac$ indexed by $t \in (0, 1)$ with the tangent field $\nabla \varphi_t$ driving its dynamics. 

\begin{restatable}[Wasserstein Parallel Transport PDE]{corollary}{} \label{Wasserstein parallel transport pde}
    A vector field $v_t$ along $\mu_t$ is parallel along $\mu_t$ with respect to $\nabla_{(\nabla\varphi_t)}^{W_2}$ if
    \begin{equation*}
        \boxed{\nabla \cdot \left(\mu_t\left(\partial_t v_t + \nabla v_t \cdot \nabla \varphi_t\right)\right) = 0 \qquad \text{for a.e. $t \in (0,1).$}}
    \end{equation*}
\end{restatable}
\noindent \textit{Proof.} Due to \Cref{parallel transport on general manifolds} and \Cref{Wasserstein covariant deriv} we know that $v_t$ is parallel along $\mu_t$ if $\Pi_{\mu_t}\left(\partial_tv_t + \nabla v_t \cdot \nabla\varphi_t\right) = 0$ for almost every $t \in (0,1).$ This is tantamount to the requirement that $\nabla \cdot \left(\mu_t(\partial_tv_t + \nabla v_t \cdot \nabla\varphi_t)\right) = 0$ since the $L_2(\mu_t)$ projection simply discards the divergence free portion of the vector field. \qed{}

The same approach is directly applicable for the Hellinger-Kantorovich case. Consider a smooth curve $\mu_t$ indexed by $t \in (0,1)$ with velocity field $(\nabla\varphi_t, \varphi_t)$. 

\begin{restatable}[Hellinger-Kantorovich Parallel Transport PDE]{corollary}{} \label{HK parallel transport pde}
    A vector field $(v_t, \beta_t)$ along $\mu_t$ is parallel along $\mu_t$ with respect to $\nabla_{(\nabla\varphi_t, \varphi_t)}^{\text{HK}}$ if
    \begin{align*}
        \boxed{\nabla \cdot (\mu_t (\partial_tv_t + \nabla v_t \cdot \nabla \varphi_t + 2\varphi_t v_t + 2g_t \nabla\varphi_t)) = 0 
        \qquad \text{and} \qquad 
        \partial_tg_t + \frac{1}{2}\langle\nabla g_t, \nabla\varphi_t\rangle + 2\varphi_t g_t = 0.}
    \end{align*}
\end{restatable}
\noindent \textit{Proof.} By \Cref{parallel transport on general manifolds} and \Cref{HK covariant deriv}, a vector field $(v_t, \beta_t)$ is parallel along $\mu_t$ if
\begin{align}
    \Pi_{\mu_t}\left(\begin{bmatrix}
        \partial_tv_t + \nabla v_t\cdot \nabla\varphi_t + 2\varphi_tv_t +  2g_t\nabla\varphi_t \\
        \partial_tg_t + \frac{1}{2}\langle\nabla g_t, \nabla\varphi_t\rangle + 2\varphi_t g_t
    \end{bmatrix}\right) = 0 \qquad \text{for a.e. $t \in (0,1).$ }
\end{align}
The orthogonal projection in the HK case is a Helmholtz decomposition of the vectorial portion of the argument (discarding the zero divergence component) and it leaves the mass creation term unchanged. Thus, we are left with the two conditions that we set out to prove. \qed{}
\\

\noindent \textbf{Approximate Parallel Transport via the Exponential Map.}
While \Cref{Wasserstein parallel transport pde} and \Cref{HK parallel transport pde} describe parallel transport, solving these PDEs in practice may be challenging. To this end, we describe an alternative approximate but tractable approach that leverages the first order properties of the exponential map. 

For a Riemannian manifold $\man$, one way to define the differential of the exponential map is as follows. Let $t \mapsto (x_t, u_t) \in T\man$ be a smooth curve through the tangent bundle, and define the induced curve $t \mapsto y_t := \exp_{x_t}(u_t).$ Then the differential of $\exp : T\man \rightarrow \man$ at the point $(x_t, u_t)$ in the direction $(\dot x_t, \nabla_{\dot x_t}u_t)$ (denoted by $d(\exp_{x_t}(u_t))(\dot x_t,  \dot u_t)$) is equal to $\dot y_t.$ In the general case of $\mathcal{P}_{2, \text{ac}}(\man)$, \citet{gigli2012second} shows that the differential of the exponential map can be derived from an ODE on the base manifold $\man$. Note that we use $\textbf{exp}$ to denote the exponential map on $\pac$ or $\posm$, while we use $\exp$ to denote the exponential map on the base space (which in both cases is $\mathbb{R}^d$). 


\begin{restatable}[Theorem 7.2, \citet{gigli2012second}]{thm}{}
    Let $\mu_t$ be a compactly-supported regular curve through $\mathcal{P}_{2, \text{ac}}(\man)$ with velocity field $\nabla \varphi_t$, and let $u_t$ be an absolutely continuous vector field along it. Define $\nu_t := \mathbf{exp}_{\mu_t}(u_t) = (\exp(u_t))_\#\mu_t$, $(\exp(u_t))(x) := \exp_x(u_t(x))$ and assume that $\exp_x(u_t(x))$ is $\mu_t$-essentially invertible for a.e. $t$. Further assume that the vectors $u_t$ are uniformly bounded in $L_{\infty}(\mu_t)$ and that 
    \[\int_0^1 \left\|\frac{d}{dt} u_t\right\|^2_{L_2(\mu_t)} dt < \infty.\]
    Then $\nu_t$ is absolutely continuous and its velocity vector field $w_t$ is given by
    \[w_t = \Pi_{\nu_t}\left(j_{\mu_t, u_t}\left(\nabla \varphi_t, \frac{d}{dt} u_t\right) \right)\]
    and for $\mu_t$-a.e. $x$, $j_{\mu_t, u_t}(a, b) (\exp_x(u_t(x))) := J_x(s)$ and $J_x$ solves
    \[\nabla_{\dot\gamma_x(s)}^2 J_x(s) + R(J_x(s), \dot \gamma_x(s))\dot \gamma_vx(s) = 0, \quad \gamma_x(s) := \exp_x(su_t(x)), \quad J_x(0) = a(x), \quad \nabla_{\dot\gamma_x(0)}J_x(0) = b(x)\]
    and $R(\cdot, \cdot)$ is the Riemann curvature tensor of $\man$.
\end{restatable}


Critically, the differential of the exponential map can be used to approximate parallel transport to the second order. We will prove this in the case of both $(\pac, W_2)$ and $(\posm, \text{HK})$. 

\begin{restatable}[Approximation of $\mathcal{P}_{2, \text{ac}}(\man)$ Parallel Transport]{thm}{}
    Let $\mu_t: [0, 1] \rightarrow \mathcal{P}_{2, \text{ac}}(\man)$ be a regular curve with velocity field $\nabla \varphi_t$ and let $\nu \in \mathcal{P}_{2, \text{ac}}(\man)$ be the target measure. Assume that there exists a $u_t$ such that $s \rightarrow \nu_s := \textbf{exp}_{\mu_t}(su_t)$, $s \in [0, 1]$ traces out a regular geodesic between $\mu_t$ and $\nu$ that has a velocity field $\nabla\psi_s$ with a Lipschitz constant uniformly bounded in $s$. Further assume that $x\mapsto\exp_x(su_t(x))$ is $\mu_t$ essentially invertible for a.e. $t$. Then it holds that 
    \[w_t^s = \Pi_{\nu_s}\left(j_{\mu_t, su_t}\left(\nabla \varphi_t, s\dot u_t\right) \right)\]
    approximates $\mathcal{T}_0^s(\nabla \varphi_t)$, the parallel transport of $\nabla\varphi_t$ along the geodesic $\textbf{exp}_{\mu_t}(su_t)$ in the sense that
    \[\|w_t^s - \mathcal{T}_0^s(\nabla\varphi_t)\|_{L_2(\nu_s)} = O(?).\]
\end{restatable}
\noindent \textit{Proof.} Let $\mathbf{T}(t, s, \cdot)$ be the flow map of the regular curve $\mu_t$ and let $u \in L_2(\mu_s).$ We'll define the following function, 
\begin{equation*}
    \tau_{t_0}^{t_1}(u)(x) := \begin{cases}
        \text{the parallel transport of $u(\mathbf{T}(t_1, t_0, x))$ on $\man$ along the curve} \\
        r \mapsto \mathbf{T}(t_1, r, x) \text{ from $r = t_0$ to $r = t_1$ }.
    \end{cases}
\end{equation*}
By the triangle inequality, we have
\begin{align*}
    \|w_t^s - \mathcal{T}_0^s(\nabla \varphi_t)\|_{L_2(\nu_s)} \leq \|w_t^s - \Pi_{\nu_s}(\tau_0^s(\nabla \varphi_t))\|_{L_2(\nu_s)} + \|\Pi_{\nu_s}(\tau_0^s(\nabla \varphi_t)) - \mathcal{T}_0^s(\nabla \varphi_t)\|_{L_2(\nu_s)}.
\end{align*}
By equation 4.18 of \citet{gigli2012second}, we have that the second term satisfies
\begin{align*}
    \|\Pi_{\nu_s}(\tau_0^s(v_t)) - \mathcal{T}_0^s(\nabla \varphi_t)\|_{L_2(\nu_s)} \leq C^2\|u_t\|_{L_2(\mu_t)}\left(\int_0^sL(\nabla\psi_r)dr\right)^2
\end{align*}
where $C < \infty$ is a constant and $L(\nabla\psi_r)$ is the largest Lipschitz constant of the velocity field generating the geodesic $\mathbf{exp}_{\mu_t}(su_t).$ Our assumption that $s \mapsto \mathbf{exp}_{\mu_t}(su_t)$ is a regular geodesic guarantees that $\int_0^sL(\nabla\psi_r)dr < \infty$, and our bounded Lipschitz constant assumption ensures that the right-hand side above is $O(s^2)$. Now we will handle the first term in the triangle inequality. Note that $\Pi_{\nu_s}: L_2(\nu_s) \rightarrow T_{\nu_s}\pac$ is a projection onto a closed and convex subset, and it is therefore non-expansive. This allows us to say
\begin{align*}
    \|w_t^s - \Pi_{\nu_s}(\tau_0^s(v_t))\|_{L_2(\nu_s)} &\leq \left\|j_{\mu_t, su_t}\left(\nabla \varphi_t, s\dot u_t\right)  - \tau_0^s(\nabla \varphi_t) \right\|_{L_2(\nu_s)}.
\end{align*}
Note that by $\mu_t$-essential invertibility of $\exp_x(su_t(x))$ we know that $j_{\mu_t, su_t}$ is well defined $\mathbf{exp}_{\mu_t}(su_t)$ a.e.. By linearity we have that $j_{\mu_t, su_t}(\nabla \varphi_t, s\dot u_t) = j_{\mu_t, su_t}(0, s\dot u_t)  + j_{\mu_t, su_t}(\nabla \varphi_t, 0)$, yielding
\begin{align*}
    \|w_t^s - \Pi_{\nu_s}(\tau_0^s(v_t))\|_{L_2(\nu_s)} &\leq \left\|j_{\mu_t, su_t}\left(\nabla \varphi_t, 0\right)  - \tau_0^s(\nabla \varphi_t) \right\|_{L_2(\nu_s)} + \|j_{\mu_t, su_t}(0, s\dot u_t) \|_{L_2(\nu_s)}.
\end{align*}

\break
\section{The cone metric}
An equivalent construction of the Wasserstein Fisher Rao space comes through lifting measures on the base space $\mathbb{R}^d$ to $\mathfrak{C}(\mathbb{R}^d)$, the metric cone of the base space ($\mathbb{R}^d$). 
\begin{restatable}[Metric cone]{defn}{}
    For a Riemannian manifold $(\man, g)$, the metric cone is
    \[\mathfrak{C}(\man) = \man \times \mathbb{R}_{+} / \{(x, 0) \sim (y, 0)\}\]
    where a point in $C(\man)$ is written as $(x, r)$. The metric tensor on the cone is given by
    \[g_{(x,r)} = dr^2 + r^2g_x.\]
\end{restatable} 
The metric (as one would have in a metric space) is given by
\begin{equation}
    d_{\mathfrak{C}}((x_0, r_0), (x_1, r_1))^2 = r_0^2 + r_1^2 - 2r_0r_1\cos(|x_0 - x_1| \land \pi).
\end{equation}
\tls{Its not clear to me that this metric follows directly the metric tensor. I believe the clipping in the argument of the cosine is an artificial addition...} \ndfg{Check pages 88 and following in Burago Burago Ivanov for exactly this!} Now, given a measure $\lambda \in M(\mathfrak{C})$ we can project it to a measure $\mathfrak{B}\lambda \in M(\mathbb{R}^d)$ with
\begin{equation}
    \int f(x)d\mathfrak{B}\lambda(x) = \int r^2 f(x) d\lambda(x, r) \qquad \text{for all test functions }f.
\end{equation}
Clearly, there are many possible ways to define a lift of a measure -- any $\lambda$ of the form
\begin{equation}
    \lambda(dx, dr) = \eta(x, dr)\mathfrak{B}\lambda(dx) \qquad \text{with } \qquad \int_0^\infty r^2\eta(x, dr) = 1 \text{ \quad for $\mathfrak{B}\lambda$ a.e. $x$}
\end{equation}
will do. To see this plug into the projection formula above and we obtain
\[\int r^2 f(x)\eta(x, dr)\mathfrak{B}\lambda(dx) = \int f(x) \left(\int r^2 \eta(x, dr) \right) \mathfrak{B}\lambda(dx) = \int f(x) \mathfrak{B}\lambda(dx) \]
as desired. \citet{liero2016optimal} show that one can characterize the Wasserstein-Fisher-Rao distance on a base space to the Wasserstein on the lifted cone space,
\begin{equation} \label{eq: cone wasserstein and hk equivalence}
    \text{HK}(\mu_0, \mu_1) = \min\left\{W_{\mathfrak{C}}(\lambda_0, \lambda_1) \,\Big|\, \mathfrak{B}\lambda_0 = \mu_0, \mathfrak{B}\lambda_1 = \mu_1\right\}.
\end{equation}
% \tls{If we want to compute HK distances we need to use the logarithmic entropy transport functional approach. I'll add a section for this soon.}

\subsection{Logarithmic-entropy-transport functional approach}
To compute the HK distance we will use a equivalent characterization that is more computationally amenable. Let $\eta_j = \Pi^j_\#\eta$ and define the Hellinger-Kantorovich entropy-transport functional as
\begin{equation}
    E(\eta; \mu_0, \mu_1) = \int_{\mathbb{R^d}}F\left(\frac{d\eta_0}{d\mu_0}\right)d\mu_0 + \int_{\mathbb{R^d}}F\left(\frac{d\eta_1}{d\mu_1}\right)d\mu_1 + \int_{\mathbb{R}^d\times \mathbb{R}^d}c(|x_0-x_1|)d\eta
\end{equation}
where
\begin{equation}
    c(L) = \begin{cases}
        -2\log(\cos(L)) & L < \pi/2 \\
        \infty & L \geq \pi/2
    \end{cases}
\end{equation}
and 
\begin{equation}
    F(\rho) = \rho\log\rho - \rho + 1.
\end{equation}
\citet{liero2016optimal} show that 
\begin{align} \label{eq: let hk equivalence}
    \text{HK}(\mu_0, \mu_1)^2 &= \inf\left\{E(\eta;\mu_0, \mu_1) \, |\, \eta  \in M(\mathbb{R}^d\times \mathbb{R}^d), \nu_j  \ll \mu_j\right\}.
\end{align}
They also show that the logarithmic-entropy-transport functional gives us a way to compute optimal lifts in the sense of \cref{eq: cone wasserstein and hk equivalence}. In particular, assume that $\eta \in M(\mathbb{R}^d\times \mathbb{R}^d)$ is a minimizer of $E(\cdot\,; \mu_0, \mu_1)$ and for the marginals $\eta_i$ consider the Lebesgue decomposition $\mu_i = \sigma_i\eta_i + \mu_i^\perp$. Then the transport plan $\gamma_{\eta} \in M(\mathfrak{C}_{\mathbb{R}^d} \times \mathfrak{C}_{\mathbb{R}^d})$ defined by
\begin{multline}
    \gamma_{\eta}(dz_0, dz_1) = \delta_{\sqrt{\sigma_0(x_0)}}(dr_0)\delta_{\sqrt{\sigma_1(x_1)}}(dr_1)\eta(dx_0, dx_1) \\
    + \delta_1(dr_0)\mu_0^\perp(dx_0)\delta_{\mathfrak{o}}(dz_1) + \delta_1(dr_1)\mu_1^\perp(dx_1)\delta_{\mathfrak{o}}(dz_0) 
\end{multline}
and the associated lifts $\lambda_i = \Pi_{\#}^i\gamma_\eta$ are optimal for \cref{eq: cone wasserstein and hk equivalence}. This means we can solve the logarithmic-entropy-transport functional approach and use the optimal plan to construct appropriate lifts of the original measures onto the cone. 

\subsection{Relevant geometric objects on the cone}

Taking $\man = \mathbb{R}^d$ to be the base space of our cone construction, we can construct tangent vectors as
\begin{equation}
    (u, a) \in T_{(x,r)}\mathfrak{C}(\mathbb{R}^d)
\end{equation}
where $u \in T_x\mathbb{R}^d \cong \mathbb{R}^d$ and $a \in \mathbb{R}$. Equivalently, it can be written as $u + a\partial_r$ where $\partial_r$ is the canonical radial vector field. The metric cone described above falls under the \say{warped product} construction, which we define explicitly below. 
\begin{restatable}[Warped product, \citet{petersen2006riemannian}]{defn}{}
    Given a Riemannian metric $(\man, g)$, a warped product (over $I$) is defined as a metric on $I \times \man$, where $I \subset \mathbb{R}$ is an open interval with metric 
    \[g = dr^2 + \rho^2(r)g\]
    where $\rho > 0$ on $I$. This is sometimes denoted $I \times_{\rho} \man$, where $I$ is called the base space and $\man$ is called the fiber.
\end{restatable}
Having established this definition, it is clear that the cone construction is a warped product using $\rho(r) = r$. With this in mind, we can used the theory of warped products to study the cone metric. Before moving on, we need to establish the \say{lift} of a vector field on a manifold to a vector field on a product manifold.

\begin{restatable}[Lifting, \citet{o1983semi}]{defn}{}
    Let $\pi: \mathcal{N} \times \man\rightarrow \mathcal{N}$ be the projection onto the first factor $\pi(p, q) = p$. 
    \begin{enumerate}
        \item If $x \in T_p\mathcal{N}$ and $q \in \mathcal{M}$ then the lift $\tilde{x}$ of $x$ to $(p,q)$ is the unique vector in $T_{(p,q)}(\mathcal{N} \times q)$ such that $d\pi(\tilde{x}) = x$. We denote the set of all such horizontal lifts $\mathfrak{L}(\mathcal{N}).$
        \item If $A \in \mathfrak{X}(\mathcal{N})$, then the lift of $A$ to $\mathcal{N} \times \man$ is the vector field $B$ whose values at each $(p,q)$ is the lift of $A_p$ to $(p,q).$
    \end{enumerate}
    Vertical lifts are defined in the same way but using the projection onto the second factor $\sigma: \mathcal{N} \times \man \rightarrow \man$. The set of all vertical lifts are denoted $\mathfrak{L}(\man).$ 
\end{restatable}



\subsubsection{The Levi-Civita connection}
Proposition 35 of \citet{o1983semi} derives the Levi-Civita connection for warped products. 

\begin{restatable}[Proposition 35, \citet{o1983semi}]{prop}{}
    Let $\man$ and $I$ be Riemannian manifolds with connections denoted $^{\man}\nabla$ and $^I\nabla$ respectively. On $I \times_\rho \man$, if $X, Y \in \mathfrak{L}(\man)$ and $V, W \in \mathfrak{L}(I)$ then
    \begin{enumerate}
        \item $\nabla_VW \in \mathfrak{L}(I)$ is the horizontal lift of $^{I}\nabla_VW$ on $I$.
        \item $\nabla_VX = \nabla_XV = \left(\frac{V(\rho)}{\rho}\right)X$.
        \item $\nabla_XY = \,^\man\nabla_XY - \left(\frac{\langle X, Y\rangle}{\rho}\right)\text{grad}(\rho)$.
    \end{enumerate}
\end{restatable}
\noindent Now we can directly apply this proposition to compute the connection on $\mathfrak{C}(\mathbb{R}^d).$ 

\begin{restatable}[Levi-Civita connection on $\mathfrak{C}(\mathbb{R}^d)$]{corollary}{} \label{corollary: connection on cone}
    Let $X, Y \in \mathfrak{L}(\mathbb{R}^d)$ and let $V, W \in \mathfrak{L}(\mathbb{R}_+)$. In particular, let $X = (X^1(x), \dots, X^n(x))$ and $Y= (Y^1(x), \dots, Y^d(x))$ in standard coordinates. Let $\nabla$ denote the Levi-Civita connection on $\mathfrak{C}(\mathbb{R}^d) \setminus \{\mathfrak{o}\}$ where $\mathfrak{o}$ is the base of the cone and let $(x,r) \in \mathfrak{C}(\mathbb{R}^d)$. It holds that 
    \begin{enumerate}
        \item $\nabla_VW$ is given by 
        \[\nabla_VW(x,r) = \left(0, V(r)\frac{\partial W(r)}{\partial r}\right).\]
        \item $\nabla_VX(x,r) = \nabla_XV(x,r) = \frac{1}{r}X(x,r)$.
        \item Finally,
        \[\nabla_XY(x,r) = \left(\sum_{i = 1}^dX^i(x) \frac{\partial Y(x)}{\partial x^i},  - \frac{\langle X(x), Y(x) \rangle_{\mathfrak{C}}}{r}\right). \]
    \end{enumerate}
\end{restatable}

\subsubsection{Geodesics}

Fortunately, geodesics on the cone space are available in closed form, and we will make them explicit here \cite{liero2016optimal}. Given two points $z_0 = (x_0, r_0)$ and $z_1 = (x_1, r_1)$, a constant-speed geodesic between them can be expressed as
\begin{equation} \label{eq: cone geodesics}
    z(t) = \left[x(t; z_0, z_1) , r(t; z_0, z_1)\right]
\end{equation}
where
\begin{equation} \label{eq: cone geodesic radial}
    r(t; z_0, z_1)^2 = (1-t)^2r_0^2 + t^2r_1^2 + 2t(1-t)r_0r_1\cos(\min\{\pi, \|x_1 - x_0\|_{2}\})
\end{equation}
and 
\begin{equation} \label{eq: cone geodesic base}
    x(t; z_0, z_1) = (1 - \rho(t; z_0, z_1))x_0 + \rho(t; z_0, z_1)x_1
\end{equation}
with
\begin{equation}
    \rho(t; z_0, z_1) = \begin{cases}
        \frac{1}{\|x_1 - x_0\|_2}\arccos\left(\frac{(1-t)r_0 + tr_1\cos|x_1 - x_0|}{r(t; z_0, z_1)}\right) & \text{for }\|x_1 - x_0\|_2 < \pi \\
        \frac{1}{2}\left(1 + \text{sign}((1-t)r_0 + tr_1)\right) & \text{otherwise}
    \end{cases}.
\end{equation}
To make this a \textit{unit} speed geodesic, one can re-parameterize with $s = d_{\mathfrak{C}}(z_0, z_1)t.$ 

Now, given an optimal coupling $\gamma$ between two (probability) measures $\lambda_1, \lambda_2 \in M(\mathfrak{C})$, we can compute the Wasserstein geodesic interpolant between them as $z(s; \cdot, \cdot)_{\#}\gamma$. One can also obtain these equations through the Levi-Civita connection defined in the previous section, as it holds that for any geodesic $\gamma$, $\nabla_{\dot \gamma}\dot \gamma = 0$. Expanding this equation using the decomposition $\dot \gamma = v + \dot r\partial_r$ \Cref{corollary: connection on cone} yields
\begin{align}
    0 = \nabla_{\dot \gamma}\dot \gamma &= \nabla_{v + \dot r \partial_r}(v + \dot r\partial_r) \\
    &= \nabla_vv + \nabla_v \dot r \partial_r + \dot r \nabla_{\partial_r}v + \dot r \nabla_{\partial_r}(\dot r \partial_r) \\
    &= \dot v - r\|v\|_2^2\partial_r + \ddot r \partial_r + \dot r \nabla_v\partial_r + \dot r \nabla_{\partial_r}v + 0\cdot  \partial_r + \dot r^2 \nabla_{\partial_r}\partial_r \\
    &= \left(\dot v + 2v\frac{\dot r}{r}\right) + \left(\ddot r  - r\|v\|_2^2\right)\partial_r.
\end{align}
Thus, the geodesic equations on the cone are
\begin{equation}\label{eq: geodesic equations on the cone}
    \dot v + 2v\frac{\dot r}{r} = 0 \qquad \text{and} \qquad \ddot r  - r\|v\|_2^2 = 0.
\end{equation}
This also has an interesting consequence. If we differentiate the following quantity against $t$
\begin{align*}
    \frac{d}{dt}(r^2v) &= 2r\dot rv + r^2 \dot v \\
    &= 2r \dot r v + r^2(-2v\dot r/r) \\
    &= 0.
\end{align*}
Thus $r^2v$ is constant in time. We will leverage this property to simplify the Jacobi equation. 

\subsubsection{The curvature tensor}

The last piece of the puzzle is the Riemann curvature tensor, as we want to be able to compute $R(u, \dot\gamma(t))\dot\gamma(t)$ to solve the Jacobi equation described in the next subsection. The curvature tensor of a (pseudo)-Riemannian manifold measures the noncommutativity of the covariant derivative, and is defined as
\begin{equation}
    R(X,Y)Z = \nabla_X \nabla_YZ - \nabla_Y\nabla_XZ - \nabla_{[X,Y]}Z
\end{equation}
where $X,Y, Z \in \mathfrak{X}(\man)$ are vector fields over the manifold. In the situation of interest, we have $\man = \mathfrak{C}(\mathbb{R}^d)$ - we will appeal to another result from the theory of warped-product manifolds to obtain an expression for the Riemann curvature tensor as a function of the curvature tensors of the base space and the fiber. 
\begin{restatable}[Proposition 7.42, \citet{o1983semi}]{prop}{} \label{prop: warped product curvature tensor}
      Let $\man$ and $I$ be Riemannian manifolds with Riemann curvature tensors denoted $^{\man}R$ and $^IR$ respectively. On $I \times_\rho \man$, if $X, Y, Z \in \mathfrak{L}(\man)$ and $U, V, W \in \mathfrak{L}(I)$ then
    \begin{enumerate}
        \item $R(U, V)W \in \mathfrak{L}(I)$ is the horizontal lift of $^IR(U,V)W$ on $I$.
        \item $R(Y,U)V = (H^\rho(U,V)/\rho)Y$ where $H^\rho$ is the Hessian of $\rho$.
        \item $R(U,V)Y = R(Y,Z)U = 0.$
        \item $R(U,Y)Z = (\langle Y,Z\rangle/\rho)\nabla_U\,\text{grad}(\rho)$\item $R(Y,Z)X = \,^\man R(Y,Z)X - (\langle \text{grad}\rho, \text{grad}\rho\rangle/\rho^2)\left\{\langle Y,X \rangle  Z - \langle Z, X\rangle Y\right\}.$
    \end{enumerate}
\end{restatable}
Now we'll drop the dependence on time and write $J = j +  b\partial_r$ and $\dot\gamma = v +  \dot r\partial_r$ where $j, v \in \mathfrak{L}(\mathbb{R}^d)$ are vertical lifts, and $b\partial_r,\dot r\partial_r \in \mathfrak{L}(I)$ are horizontal lifts. By linearity
\begin{multline}
    R(J, \dot\gamma)\dot\gamma = R(j, v)v + \dot rR(j, \partial_r)v + bR(\partial_r, v)v + b\dot rR(\partial_r, \partial_r)v \\
    + \dot rR(j, v)\partial_r + \dot r^2R(j, \partial_r)\partial_r + \dot rbR(\partial_r, v)\partial_r + b\dot r^2R(\partial_r, \partial_r)\partial_r.
\end{multline}
\Cref{prop: warped product curvature tensor} indicates that any term $R(a,b)c$ where $a,b \in \mathfrak{L}(\mathbb{R}^d)$ and $c \in \mathfrak{L}(I)$ (or vice versa) is zero. Also note that $I$ is flat, so $^IR \equiv0$, implying $r^2R(\partial_r, \partial_r)\partial_r = 0$. Thus,
\begin{align}
    R(J, \dot\gamma)\dot\gamma = R(j, v)v + \dot rR(j, \partial_r)v + bR(\partial_r, v)v + \dot r^2R(j, \partial_r)\partial_r + \dot rbR(\partial_r, v)\partial_r.
\end{align}
\Cref{prop: warped product curvature tensor} indicates that 
\begin{align}
    \dot rR(j, \partial_r)v &= -\dot rR(\partial_r, j)v & \text{by antisymmetry}\\
    &= \left(\frac{\langle j, v\rangle_{\mathbb{R}^d}}{\rho}\right)\nabla_{\partial_r}\text{grad}(\rho) \\
    &= \left(\frac{\langle j, v\rangle_{\mathbb{R}^d}}{\rho}\right)\nabla_{\partial_r}\partial_r \\
    &= 0
\end{align}
and $bR(\partial_r, v)v = 0$ since $\rho(r) = r$. We also have
\begin{align}
    \dot r^2R(j, \partial_r)\partial_r &= \left(\frac{H^{\rho}(\partial_r, \partial_r)}{\rho}\right)j \\
    &= 0
\end{align}
since $H^{\rho}\equiv 0$ and $\dot rbR(\partial_r, v)\partial_r = 0$ by the same logic. Putting it together, we are left with
\begin{align}
    R(J, \dot\gamma)\dot\gamma &= R(j,v)v \\
     &= \,^{\mathbb{R}^d}R(j,v)v - \frac{\langle \text{grad} \rho, \text{grad}\rho\rangle}{\rho^2}\left(\langle j, v \rangle v -  \langle v, v\rangle j\right) \\
     &= -\frac{1}{r^2}\left(\langle j, v\rangle_{\mathfrak{C}} \cdot v - \|v\|^2_{\mathfrak{C}}\cdot j\right) \\
     &= j\|v\|^2_{2} - v\langle j,v\rangle_{\mathbb{R}^d}. \label{eq: jacobi riemann expression}
\end{align}
Now we can plug this directly into the Jacobi equation. Note that this scalar quantity above should really be lifted to the cone (and we will do so in the subsequent derivations).

\subsubsection{Putting it together: Jacobi fields}

To compute Jacobi fields on the cone space described above given initial conditions $J(0)$ and $J'(0)$, one needs to solve the Jacobi equation
\begin{equation}
    \nabla_{\dot\gamma}^2J(t) + R(J(t), \dot{\gamma}(t))\dot{\gamma}(t) = 0
\end{equation}
where $\gamma$ is the geodesic of interest on the cone space and $R$ is the Riemann curvature tensor. Let $J = j + b\partial_r$ and let $\dot\gamma = v + \dot r \partial_r$. Plugging in the result from \cref{eq: jacobi riemann expression} yields
\begin{align}
    \nabla_{v + \dot r \partial_r}^2(j + b\partial_r) + j\|v\|^2_{2} - v\langle j,v\rangle_{\mathbb{R}^d} = 0.
\end{align}
Now we'll expand the first term as follows
\begin{align}
    \nabla_{v + \dot r \partial_r}(j + b\partial_r) &= \nabla_vj + \nabla_v(b\partial_r) + \nabla_{\dot r\partial_r}j + \nabla_{\dot r\partial_r}(b\partial_r) \\
    &= \nabla_vj + v(b)\partial_r + b\nabla_v\partial_r + \dot r\nabla_{\partial_r}j + \nabla_{\dot r\partial_r}\partial_r + 0.
\end{align}
Continuing, we have
\begin{align}
    \nabla_{v + \dot r \partial_r}(j + b\partial_r)&= \,^{\mathbb{R}^d}\nabla_vj - \frac{\langle v, j\rangle_{\mathfrak{C}}}{\rho}\text{grad}\rho + \dot b\partial_r +  \left(\frac{\partial_r(\rho)}{\rho}\right)bv  + \frac{\dot r\partial_r(\rho)}{\rho}j & \text{ by \Cref{corollary: connection on cone}}\\
    &= \left(\dot j + \frac{v}{r}b  + \frac{\dot r}{r}j, \dot b - \frac{\langle v, j\rangle_{\mathfrak{C}}}{r}\right) \\
    &= \left(\dot j + \frac{v}{r}b  + \frac{\dot r}{r}j, \dot b - r\langle v, j\rangle_{\mathbb{R}^d}\right)
\end{align}
For convenience, let $A = \dot j + \frac{v}{r}b  + \frac{\dot r}{r}j$ and let $c = \dot b - r\langle v, j\rangle_{\mathbb{R}^d}$. Now we'll apply the covariant derivative again,
\begin{align}
    \nabla^2_{\dot\gamma}J(t) &= \nabla_{\dot\gamma}(A + c\partial_r) \\
    &= \nabla_v A + \dot r\nabla_{\partial_r}A + \dot c\partial_r + c\nabla_{\dot\gamma}\partial_r\\
    &= \dot A  - \frac{\langle v, A\rangle_{\mathfrak{C}}}{r}\partial_r + \dot r\nabla_{\partial_r}A + \dot c\partial_r + c\nabla_{v}\partial_r + \dot rc\nabla_{\partial_r}\partial_r \\
    &= \dot A  - r\langle v, A\rangle_{\mathbb{R}^d}\partial_r + \frac{\dot r}{r}A + \dot c\partial_r + \frac{cv}{r} \\
    &= \left(\dot A  + \frac{\dot r}{r}A + \frac{c}{r}v\right) + \left(\dot c - r\langle v, A \rangle_{\mathbb{R^d}}\right)\partial_r.
\end{align}
Plugging back into the Jacobi equation yields
\begin{equation}
    \left(\dot A  + \frac{\dot r}{r}A + \frac{c}{r}v + j\|v\|_2^2 - v\langle j, v\rangle_{\mathbb{R^d}}\right) + \left(\dot c - r\langle v, A \rangle_{\mathbb{R^d}}\right)\partial_r = 0.
\end{equation}
Expanding further, we have
\begin{align}
    \dot A &= \ddot j + \frac{r\dot v - v \dot r}{r^2}b + \frac{v}{r}\dot b + j\left(\frac{r\ddot r - \dot r ^2}{r^2}\right) + \frac{\dot r}{r}\dot j\\
    &= \ddot j + \left(\frac{\ddot r}{r} - \frac{\dot r^2}{r^2}\right)j + \frac{\dot r}{r}\dot j + \left(\frac{\dot b}{r} - \frac{b \dot r}{r^2}\right)v + \frac{b}{r}\dot v
\end{align}
and 
\begin{align}
    \dot c &= \ddot b - \dot r \langle v, j\rangle - r\langle \dot v, j \rangle - r\langle v, \dot j\rangle.
\end{align}
The geodesic equations in \cref{eq: geodesic equations on the cone} indicate that $\dot v = -2v\dot r/r$ and $\ddot r = r\|v\|_2^2$. Substituting this in yields
\begin{align*}
    \dot A = \ddot j + \left(\|v\|_2^2 - \frac{\dot r^2}{r^2}\right)j + \frac{\dot r}{r}\dot j + \left(\frac{\dot b}{r} - \frac{b \dot r}{r^2}\right)v  -2v\frac{b\dot r}{r^2}.
\end{align*}
Adding the $A\dot r/r$ term from the Jacobi equation yields
\begin{align*}
    \dot A + \frac{\dot r}{r}A = \ddot j +  2\frac{\dot r}{r}\dot j +\|v\|_2^2j + \left(\frac{\dot b}{r} - \frac{2b \dot r}{r^2}\right)v
\end{align*}
and plugging back into the base component of the Jacobi equation
\begin{align}
    0 &= \ddot j +  2\frac{\dot r}{r}\dot j +\|v\|_2^2j + \left(\frac{\dot b}{r} - \frac{2b \dot r}{r^2}\right)v + \frac{c}{r}v + j\|v\|_2^2 - v\langle j, v\rangle \\
    &= \ddot j +  2\frac{\dot r}{r}\dot j +\|v\|_2^2j + \left(\frac{\dot b}{r} - \frac{2b \dot r}{r^2}\right)v + \left(\frac{\dot b}{r} - \langle v, j\rangle \right)v + j\|v\|_2^2 - v\langle j, v\rangle\\
    &= \ddot j +  2\frac{\dot r}{r}\dot j +2\|v\|_2^2j + \left(\frac{2\dot b}{r} - \frac{2b \dot r}{r^2}\right)v - 2v\langle j, v\rangle \\
    &= \ddot j +  2\frac{\dot r}{r}\dot j +2\|v\|_2^2j + \frac{2}{r}\left(\dot b - b\frac{ \dot r}{r}\right)v - 2v\langle j, v\rangle.
\end{align}
Now we'll do the same for the radial component,
\begin{align*}
    \dot c - r\langle v, A\rangle &= \dot c - r\langle v, \dot j\rangle - b\|v\|_2^2 - \dot r\langle v, j\rangle \\
    &= \ddot b - \dot r \langle v, j\rangle - r\langle \dot v, j \rangle - r\langle v, \dot j\rangle  - r\langle v, \dot j\rangle - b\|v\|_2^2 - \dot r\langle v, j\rangle \\
    &= \ddot b - 2r\langle v, \dot j\rangle - b \|v\|_2^2.
\end{align*}
Thus, our Jacobi equations are

\begin{equation}
    \boxed{\ddot j +  2\frac{\dot r}{r}\dot j +2\|v\|_2^2j + \frac{2}{r}\left(\dot b - b\frac{ \dot r}{r}\right)v - 2v\langle j, v\rangle = 0}
\end{equation}
and 
\begin{equation}
    \boxed{\ddot b - 2r\langle v, \dot j\rangle - b \|v\|_2^2 = 0.}
\end{equation}


\subsubsection{Solving the Jacobi equations}
Since $r^2v$ is constant in time we will define it to be $q := r^2v$ and substitute in $v = q/r^2$, 
\begin{equation}
    \ddot j +  2\frac{\dot r}{r}\dot j +\frac{2}{r^4}\|q\|_2^2j + \frac{2}{r^3}\left(\dot b - b\frac{ \dot r}{r}\right)q - \frac{2}{r^4}q\langle j, q\rangle = 0
\end{equation}
and 
\begin{equation}
    \ddot b - \frac{2}{r}\langle q, \dot j\rangle - \frac{b}{r^4} \|q\|_2^2 = 0.
\end{equation}
Now we will decompose $j$ into $j_{\perp} + aq$, where $j_{\perp} \perp q$. It follows that $\langle j, q\rangle = a\|q\|_2^2$, $\dot j = \dot j_{\perp} + \dot aq$, and $\langle \dot j, q\rangle = \dot a \|q\|_2^2$. Thus,
\begin{equation}
    \ddot j +  2\frac{\dot r}{r}\dot j +\frac{2}{r^4}\|q\|_2^2j_{\perp} + \frac{2}{r^3}\left(\dot b - b\frac{ \dot r}{r}\right)q = 0
\end{equation}
and 
\begin{equation}
    \ddot b - \frac{2\dot a}{r}\|q\|_2^2 - \frac{b}{r^4} \|q\|_2^2 = 0.
\end{equation}
Since the LHS of the equations are equal to zero, we can project onto the orthogonal complement of $q$ and the span of $q$ to get an equivalent form for the first equation,
\begin{equation}
    \ddot j_{\perp} +  2\frac{\dot r}{r}\dot j_{\perp} +\frac{2}{r^4}\|q\|_2^2j_{\perp} = 0
\end{equation}
and 
\begin{equation}
    \ddot a + 2\frac{\dot r}{r}\dot a + \frac{2}{r^3}\left(\dot b - b\frac{ \dot r}{r}\right) = 0.
\end{equation}
Now let $u_\perp = rj_\perp$ and let $u = b/r$. By \Cref{lemma: orthog jacobi equation} we have
\begin{equation} 
    \ddot u_{\perp} + \frac{\|q\|_2^2}{r^4}u_\perp = 0.
\end{equation}
This gives us a way to solve for $u_\perp$. By \Cref{lemma: parallel jacobi equation} we have
\begin{equation} \label{eq: p and w ode}
    \dot p + 2\dot u = 0 \qquad \text{and} \qquad \dot w - 2\frac{\|q\|_2^2}{r^2}p = 0
\end{equation}
where $p = r^2\dot a$ and $w = r^2\dot u$. To solve the equations in \cref{eq: p and w ode} observe that $\dot p = (-2w/r^2)$ and $\dot w = 2p\|q\|_2^2/r^2$. Thus we have the system
\begin{align*}
    \begin{bmatrix}
        \dot p \\ \dot w
    \end{bmatrix} &= 
    \underbrace{
    \begin{bmatrix}
        0 & \frac{-2}{r(t)^2} \\
        2\frac{\|q\|_2^2}{r(t)^2} & 0 
    \end{bmatrix}}_{A(t)}
    \begin{bmatrix}
        p \\
        w
    \end{bmatrix} \\
    &= \frac{1}{r(t)^2}
    \begin{bmatrix}
        0 & -2  \\
        2\|q\|_2^2 & 0 
    \end{bmatrix}
    \begin{bmatrix}
        p \\
        w
    \end{bmatrix}.
\end{align*}
Now we'll define $\tau(t) := \int_0^tds/r^2(s)$, which means 
\[\frac{d}{dt} = \frac{1}{r(t)^2}\frac{d}{d\tau}.\]
Under this change of variables, we obtain
\begin{align*}
    \frac{d}{d\tau}
    \begin{bmatrix}
        p \\ w
    \end{bmatrix} 
    &= \underbrace{\begin{bmatrix}
        0 & -2  \\
        2\|q\|_2^2 & 0 
    \end{bmatrix}}_{M}
    \begin{bmatrix}
        p \\
        w
    \end{bmatrix}.
\end{align*}
The eigenvalues of $M$ are $\lambda = \pm 2ci$ with eigenvectors $\begin{bmatrix}
    \pm i/c & 1
\end{bmatrix}^T$ with $c = \|q\|_2$. Thus the general solution is of the form 
\begin{align*}
    \begin{bmatrix}
        p \\ w
    \end{bmatrix} &= k_1\begin{bmatrix}
        -\frac{1}{c}\sin(2c\tau) \\ \cos(2c\tau)
    \end{bmatrix} + k_2\begin{bmatrix}
        \frac{1}{c}\cos(2c\tau) \\ \sin(2c\tau)
    \end{bmatrix}
\end{align*}
Solving for initial conditions $p_0$ and $w_0$ yields
\begin{align*}
    p(t) &= -\frac{w_0}{\|q\|_2}\sin(2\|q\|_2\tau(t)) + p_0\cos(2\|q\|_2\tau(t)) \\
    w(t) &= w_0\cos(2\|q\|_2\tau(t)) + p_0\|q\|_2\sin(2\|q\|_2\tau(t)).
\end{align*}
Plugging in and integrating yields
\begin{align*}
    a(t) &= a_0 -\frac{w_0}{\|q\|_2} \int_0^t\frac{\sin(2\|q\|_2\tau(s))}{r(s)^2}ds + p_0 \int_0^t \frac{\cos(2\|q\|_2\tau(s))}{r(s)^2} ds \\
    &= a_0 + \frac{w_0}{2\|q\|_2^2}(\cos(2\|q\|_2\tau(t)) - 1) + \frac{p_0}{2\|q\|_2^2}\sin(2\|q\|_2\tau(t)).
\end{align*}
Also note that we have $p_0 = r_0^2\dot a_0$ and $w_0 = r_0^2\dot u_0 = r_0^2 (\frac{\dot b_0}{r_0} - b_0\frac{\dot r_0}{r_0^2}) = \dot b_0r_0 - b_0\dot r_0.$ Putting it together, we have
\begin{equation}
    \boxed{a(t) = a_0 + \frac{\dot b_0r_0 - b_0\dot r_0}{2\|q\|_2^2}(\cos(2\|q\|_2\tau(t))-1) + \frac{r_0^2\dot a_0}{2\|q\|_2^2}\sin(2\|q\|_2\tau(t)).}
\end{equation}
Doing the same for $u$ yields
\begin{align*}
    u(t) &= u_0 + w_0 \int_0^t\frac{\cos(2\|q\|_2\tau(s))}{r(s)^2}ds + p_0 \|q\|_2 \int_0^t \frac{\sin(2\|q\|_2\tau(s))}{r(s)^2} ds \\
    &= u_0 + \frac{w_0}{2\|q\|_2}\sin(2\|q\|_2\tau(s)) - \frac{p_0}{2} (\cos(2\|q\|_2\tau(s))-1).
\end{align*}
This implies
\begin{equation}
    \boxed{b(t) = r(t)\left[\frac{b_0}{r_0} + \frac{\dot b_0r_0 - b_0\dot r_0}{2\|q\|_2}\sin(2\|q\|_2\tau(t)) + \frac{r_0^2\dot a_0}{2} (1- \cos(2\|q\|_2\tau(t)))\right].}
\end{equation}
\tls{one can solve for $a(0)$ and $\dot a(0)$ from ivs of the Jacobi field. Will describe in detail later.} All that remains is to solve for $u_\perp$. 
Fortunately, $\ddot u_{\perp} + \frac{\|q\|_2^2}{r^4}u_\perp = 0$ falls under the umbrella of a well studied variety of ODEs called time-varying harmonic oscillator. In this case, because $r$ is not a constant function of time, this is a harmonic oscillator with varying frequency. The solution is given by \Cref{lemma: time varying oscillator}, and has the form
\begin{equation}
    u_{\perp}(t) = r(t)\left[\frac{u_{\perp,0}}{r_0}\cos(\|q\|_2\sqrt{2}\tau(t)) + \frac{\dot u_{\perp,0} r_0 - u_{\perp,0} \dot r_0}{\|q\|_2\sqrt{2}}\sin(\|q\|_2\sqrt{2}\tau(t))\right].
\end{equation}
which means 
\begin{equation}
    \boxed{j_{\perp}(t) = \frac{r_0j_{\perp,0}}{r_0}\cos(\|q\|_2\sqrt{2}\tau(t)) + \frac{r_0^2\dot j_{\perp, 0}}{\|q\|_2\sqrt{2}}\sin(\|q\|_2\sqrt{2}\tau(t)).}
\end{equation}


\begin{algorithm}[H]
    \begin{algorithmic}[1]
        \caption{Jacobi field on $\mathfrak{C}(\mathbb{R}^d)$}
        \label{alg: jacobi field}
        \Require Initial conditions $J(0) = j_0 + b_0\partial_r$ and $J'(0) = \dot j_0 + \dot b_0 \partial_r$, cone geodesic $\gamma(s) = x(s) + r(s)\partial_r$, terminal time $t$.\\
        Compute geodesic invariant $q_\gamma = \dot x(0)r(0)^2$ and pseudotime $\tau(t) = \int_0^tds/r^2(s)$. \\
        Decompose $j_0 = j_{\perp, 0} + a_0q_{\gamma}$ where $j_{\perp, 0} \perp q_{\gamma}$ and $\dot j_0 = \dot j_{\perp, 0} + \dot a_0q_{\gamma}$ where $\dot j_{\perp, 0} \perp q_{\gamma}$.\\
        Compute 
        \begin{align*}
        j_{\perp}(t) &= \frac{r_0j_{\perp,0}}{r_0}\cos(\|q_{\gamma}\|_2\sqrt{2}\tau(t)) + \frac{r_0^2\dot j_{\perp, 0}}{\|q_{\gamma}\|_2\sqrt{2}}\sin(\|q_{\gamma}\|_2\sqrt{2}\tau(t)) \\
        b(t) &= r(t)\left[\frac{b_0}{r_0} + \frac{\dot b_0r_0 - b_0\dot r_0}{2\|q_{\gamma}\|_2}\sin(2\|q_{\gamma}\|_2\tau(t)) + \frac{r_0^2\dot a_0}{2} (1- \cos(2\|q_{\gamma}\|_2\tau(t)))\right] \\
        a(t) &= a_0 + \frac{\dot b_0r_0 - b_0\dot r_0}{2\|q_{\gamma}\|_2^2}(\cos(2\|q_{\gamma}\|_2\tau(t))-1) + \frac{r_0^2\dot a_0}{2\|q_{\gamma}\|_2^2}\sin(2\|q_{\gamma}\|_2\tau(t)).
        \end{align*}
        \\ \Return $J(t) = j_{\perp}(t) + a(t)q_{\gamma} + b(t)\partial_r$.
    \end{algorithmic}
\end{algorithm}

\section{Parallel transport via Jacobi fields}

\subsection{Jacobi Fields}
Jacobi fields allow us to measure the effect of curvature on a one parameter family of geodesics. A \textit{Jacobi field} is a vector field along a geodesic $\gamma$ capturing the difference between a geodesic and a perturbed geodesic along the parameter of the family. Concretely, let $\gamma_{\tau}$ be a smooth, one-parameter family of geodesics with $\gamma_0 = \gamma$. Then
\begin{equation}
    J(t) = \frac{\partial}{\partial \tau}\gamma_{\tau}(t)\bigg|_{\tau = 0}
\end{equation}
is a Jacobi field, and it describes the infinitesimal behavior of the geodesic family about $\tau = 0$. Jacobi fields satisfy the so-called Jacobi equation
\begin{equation}
    \mathbf{D}_t^2J(t) + R(J(t), \dot{\gamma}(t))\dot{\gamma}(t) = 0
\end{equation}
where $\mathbf{D}_t$ is the covariant derivative (typically with respect to the Levi-Civita connection) $\dot{\gamma}(t) = d\gamma(t)/dt$, and $R(\cdot, \cdot)$ is the Riemann curvature tensor. A key fact about Jacobi fields is the following, which arises from the properties of the second-order ODE above. 

\begin{restatable}[Uniqueness of Jacobi Fields]{prop}{}
The value of $J(t)$ and $\mathbf{D}_tJ(t)$ at a single point along $\gamma$ completely and uniquely determines the Jacobi field along the entire geodesic.    
\end{restatable}

Observe that this result indicates that there are two ways for one to define Jacobi fields on a manifold. To establish the connection to parallel transport, we will construct a specific Jacobi field as is done by \citet{louis2018fanning}. Consider a geodesic $\gamma(t)$ connecting a point $p = \gamma(0)$ and some $q$, and let $w \in T_p\man$. Now define the $1$-parameter family of geodesics 
\[\exp_{\gamma(t)}(s(\dot{\gamma}(t) + \epsilon w))\]
parameterized by $\epsilon$. As the name \say{fanning} suggests, this family represents a collection of geodesics fanning out from the point $\gamma(t)$ -- a larger $\epsilon$ means the initial velocity is pulled further towards $w$ rather than the original geodesic velocity $\dot{\gamma}(t).$ This parameterized family of geodesics then induces the following Jacobi field,
\begin{align}
    J^{w}_{\gamma(t)}(s) = \frac{\partial}{\partial\epsilon}\bigg|_{\epsilon = 0}\exp_{\gamma(t)}(s(\dot{\gamma}(t) + \epsilon w)) \in T_{\gamma(t+s)}\man.
\end{align}
As shown by \citet{louis2018fanning} it holds that 
\begin{equation}
    P_{0, s}(w) = \frac{J^w_{\gamma(0)}(s)}{s} + O(s^2)
\end{equation}
where $P_{0, s}$ is the vector field obtained by parallel transporting $w$ along the geodesic $\gamma$. In practice, if one can compute the exponential and log maps, then the Jacobi field can be approximated by finite differences,
\[J^w_{\gamma(t)}(s) \cong \frac{\exp_{\gamma(t)}(s(\dot{\gamma}(t) + \epsilon w)) - \exp_{\gamma(t)}(s\dot{\gamma}(t))}{\epsilon}.\]
Unfortunately this construction is not directly applicable to our scenario because neither $\mathcal{P}_2(\mathbb{R}^d)$ or $\text{HK}(\mathbb{R}^d)$ are truly Riemannian manifolds. 

\subsection{A related approach}

We can pursue a related approach inspired by \citet{gigli2012second}. On a manifold $\man$, one defines the differential of the exponential map as follows. Let $t \mapsto (x_t, u_t)$ be a smooth curve through the tangent bundle of $\man$, and define the curve $y_t = \exp_{x_t}(u_t)$. Then we define the differential of the exponential map to be
\[d(\exp_{x_t}(u_t))(\dot x_t, \nabla_{\dot x_t}u_t) = \dot y_t.\]


Let $\gamma(t)$ be a geodesic in $\man$, and let $w_t$ be a vector field along it with $w_0 = w$. It holds that $w_t$ is the parallel transport of $w$ if
\[w_t = d(\exp_{\gamma(0)}(tv))(w)\]
where $v = \dot{\gamma}(0)$ and $d(\cdot)$ is the differential operator \tls{I think this is true up to rescaling... I need to prove it.} Thus, this gives us a different avenue for pursuing a tractable parallel transport approach along geodesics. If we can compute the differential of the exponential map then we can compute parallel transport.

We'll start with the case of $\mathcal{P}_2(\mathbb{R}^d)$. Given $\mu \in \mathcal{P}_2(\mathbb{R}^d)$ and $u \in L_2(\mu)$, $u:\mathbb{R}^d \rightarrow \mathbb{R}^d$ the exponential map is defined as 
\begin{equation}
    \mathbf{exp}_\mu(u) = (\exp(u))_\#\mu = (\text{Id} + u)_\#\mu
\end{equation}
where $\#$ is the pushforward operator and $\exp$ is the Euclidean exponential map. Note that I will sometimes denote the Wasserstein exponential as $\mathbf{exp}_{\mu, u}$ for convinence. Now pick a regular curve $\mu_t$ with velocity field $v_t$ and an absolutely continuous vector field $u_t$ along it. Further define the curve $t \mapsto \nu_t = \mathbf{exp}_{\mu_t}(u_t)$. \textit{Assuming $\nu_t$ is absolutely continuous}, we'll consider its velocity vector field $w_t$. One should then expect that the differential of the exponential map at the point $\mu_t, u_t$ in the direction $v_t, \dot{u}_t$ is given by $w_t$. We will eventually define the map 
\[j_{\mu, u}: T_u \mathcal{P}_2(\mathbb{R}^d) \times T_u \mathcal{P}_2(\mathbb{R}^d) \rightarrow T_{\exp_{\mu}(u)} \mathcal{P}_2(\mathbb{R}^d) \]
which will represent the \textit{differential} of the exponential map,
\begin{equation}
    j_{\mu, u}(v, w) = d(\mathbf{exp}_{\mu, u}(v, w)).
\end{equation}
Note that the differential of the exponential map takes two arguments, representing instantaneous changes in the measure $\mu$ and the tangent vector $u$. Now, define $\mathbf{T}(t, s, \cdot)$ to be the flow maps of $\mu_t$ and fix some $\varphi \in C^{\infty}(\mathbb{R}^d)$. We'll consider the derivative $\frac{d}{dt}\int\varphi d\nu_t$ shortly, but first observe that
\begin{align}
    \nu_t = \mathbf{exp}_{\mu_t}(u_t) = (\exp(u_t))_\#\mu_t = (\underbrace{\exp(u_t)}_{\text{exponential map to obtain $\nu_t$}} \circ \underbrace{\mathbf{T}(0,t, \cdot)}_{\text{flow map of $\mu_t$}})_\#\mu_0.
\end{align}
Now we can take the derivative,
\begin{align}
    \frac{d}{dt}\int\varphi d\nu_t &= \int \frac{d}{dt}\left[\varphi\circ\exp(u_t) \circ \mathbf{T}(0, t, \cdot)\right] d\mu_0 \\
    &= \int \left\langle (\nabla\varphi) \circ \exp(u_t) \circ \mathbf{T}(0, t, \cdot), \frac{d}{dt} [\exp(u_t) \circ \mathbf{T}(0, t, \cdot)] \right\rangle  d\mu_0.
\end{align} 
To obtain an expression for the second argument of the inner product, we will note the following. The map $s \mapsto \exp_{x_t}(su_t)$ defines a geodesic along the base manifold, and we are differentiating along the parameter $t$. This means that the second argument in the inner product is exactly a Jacobi field along $s \mapsto \exp_{x_t}(su_t)$ evaluated at $s = 1$. The initial conditions are 
\begin{align*}
    J_t(0) &= \frac{d}{dt}[\exp(0) \circ \mathbf{T}(0, t, \cdot)] = \frac{d}{dt}\mathbf{T}(0, t, \cdot) = v_t
\end{align*}
and
\begin{align*}
    J_t'(0) &= \frac{d}{ds} \frac{d}{dt} [\exp(su_t) \circ \mathbf{T}(0, t, \cdot)]\Big|_{s = 0} \\
    &=  \frac{d}{dt} \left[\frac{d}{ds}\exp(su_t)\Big|_{s = 0} \circ \mathbf{T}(0, t, \cdot)\right] \\
    &= \frac{d}{dt} \left[u_t \circ \mathbf{T}(0, t, \cdot)\right] \\
    &= \dot u_t  \tag{ by Proposition 3.13 of \citet{gigli2012second}.} 
\end{align*}

Note that this is a Jacobi field along the \textit{base} manifold, which in the Wasserstein case is $\mathbb{R}^d$. This means we can solve for it quite easily with the Jacobi equation: let $\gamma(s) = x + s\xi$ be a geodesic with velocity $\xi$. Any Jacobi field along $\gamma$ satisfies 
\[\frac{d^2}{ds^2}J_t(s) + R(J_t(s), \xi)\xi = 0.\]
Since $\mathbb{R}^d$ is a flat space, its Riemann curvature tensor is identically zero. Thus, the Jacobi equation is satisfied if $J''_t \equiv 0$. The desired initial conditions imply
\[J_t(s) = J_t(0) + sJ'_t(0).\]
For generality lets define the \textit{vector field} $j_{\mu, u}(\tilde{u}^1, \tilde{u}^2)(\exp_x(u(x))) = J_t(1)$, which in the Euclidean case is $\tilde{u}^1(x) + \tilde{u}^2(x)$. It follows that
\begin{align}
    \frac{d}{dt}\int \varphi d\nu_t &= \int \bigg\langle \nabla \varphi, j_{\mu, u}(v_t, \dot{u}_t)\bigg\rangle\, d\nu_t \qquad \text{for all } \varphi \in C^{\infty}(\mathbb{R}^d).
\end{align}
This means that the vector field $j_{\mu, u}(v_t, \dot{u}_t)$ satisfies the continuity equation (in the weak sense) for the absolutely continuous curve $\nu_t$. Thus,
\begin{equation}
    w_t = \Pi_{\nu_t}\left(j_{\mu, u}(v_t, \dot{u}_t)\right).
\end{equation}
Specializing this to the Euclidean case yields
\begin{equation}
    w_t = \Pi_{\nu_t}\left(v_t + \dot{u}_t\right).
\end{equation}
In the case of HK, we can instead take the base space to be $\mathfrak{C}(\mathbb{R}^d)$ and compute the Jacobi field using \Cref{alg: jacobi field}.

\section{Counterfactual Dynamics estimation}



\appendix

\section{Proofs}

\section{Additional results}

\begin{restatable}{lemma}{} \label{lemma: orthog jacobi equation}
    For the orthogonal Jacobi equation
    \begin{equation}
        \ddot j_{\perp} +  2\frac{\dot r}{r}\dot j_{\perp} + \frac{2}{r^4}\|q\|_2^2j_{\perp} = 0
    \end{equation}
    we have an equivalent characterization through
    \begin{equation}
        \ddot u_{\perp} + \frac{\|q\|_2^2}{r^4}u_\perp = 0
    \end{equation}
    with $u_\perp = rj_{\perp}$.
\end{restatable}
\noindent \textit{Proof.} We have
\begin{align}
    \dot j_{\perp} &= \frac{r\dot u_{\perp} - u_\perp \dot r}{r^2}
\end{align}
and 
\begin{align}
    \ddot j_{\perp} &= \frac{r^2(\dot r\dot u_{\perp} + r \ddot u_{\perp} - \dot u_\perp \dot r - u_\perp \ddot r) - (r\dot u_{\perp} - u_\perp \dot r)(2r\dot r)}{r^4} \\
    &= \frac{r^2(r \ddot u_{\perp} - u_\perp \ddot r) - (r\dot u_{\perp} - u_\perp \dot r)(2r\dot r)}{r^4} \\
    &= \frac{r^3 \ddot u_{\perp} - r^2u_\perp \ddot r - 2r^2\dot r\dot u_{\perp} + 2ru_\perp \dot r^2}{r^4} \\
    &= \frac{r^2 \ddot u_{\perp} - 2r\dot r\dot u_{\perp}- u_\perp\left(r \ddot r - 2 \dot r^2\right)}{r^3}.
\end{align}
Plugging back into the Jacobi equation yields
\begin{align}
    0 &= \frac{r^2 \ddot u_{\perp} - 2r\dot r\dot u_{\perp}- u_\perp r\ddot r + 2 u_{\perp}\dot r^2}{r^3} + \frac{2\dot rr\dot u_{\perp} - 2u_\perp \dot r^2}{r^3} + \frac{2}{r^5}\|q\|_2^2u_{\perp} \\
    &= \frac{r^2 \ddot u_{\perp} - u_\perp r\ddot r}{r^3} + \frac{2}{r^5}\|q\|_2^2u_{\perp}.
\end{align}
Multiplying through by $r^3$ yields $r^2 \ddot u_{\perp} - u_\perp r\ddot r + \frac{2}{r^2}\|q\|_2^2u_{\perp} = 0$ and dividing by $r^2$ yields $\ddot u_{\perp} - u_\perp \frac{\ddot r}{r} + \frac{2}{r^4}\|q\|_2^2u_{\perp} = 0$. From the geodesic equations we know that $\ddot r = r\|v\|_2^2 = \|q\|^2/r^3$. Thus,
\begin{align}
    \ddot u_{\perp} - u_\perp \frac{\ddot r}{r} + \frac{2}{r^4}\|q\|_2^2u_{\perp} &= \ddot u_{\perp} -  \frac{1}{r^4}\|q\|_2^2u_\perp + \frac{2}{r^4}\|q\|_2^2u_{\perp} \\
    &= \ddot u_{\perp} + \frac{1}{r^4}\|q\|_2^2u_\perp.
\end{align}

\qed{}

\begin{restatable}{lemma}{} \label{lemma: parallel jacobi equation}
    Let $u = b/r$ and suppose that
    \begin{equation}
        \ddot a + 2\frac{\dot r}{r}\dot a + \frac{2}{r^3}\left(\dot b - b\frac{ \dot r}{r}\right) = 0.
    \end{equation}
    and
    \begin{equation}
        \ddot b - \frac{2\dot a}{r}\|q\|_2^2 - \frac{b}{r^4} \|q\|_2^2 = 0.
    \end{equation}
    It follows that
    \begin{equation}
        \dot p + 2\dot u = 0
    \end{equation}
    and 
    \begin{equation}
        (r^2\dot u)' - 2\frac{\|q\|_2^2}{r^2}p = 0
    \end{equation}
    where $u = b/r$ and $p = r^2\dot a$.
\end{restatable}
\noindent \textit{Proof.} The first equation implies
\begin{align}
    0 &= \underbrace{(r^2\ddot a + 2r\dot r\dot a)}_{ = (r^2\dot a)'} + \frac{2}{r}\underbrace{\left(\dot b - b\frac{ \dot r}{r}\right)}_{ = r\dot u} \\
    &= (r^2\dot a)' + 2\dot u.
\end{align}
Taking $p = r^2 \dot a$ finishes the proof of the first equation. The second given equation can be rewritten as 
\begin{align}
    \ddot b  - \frac{b\ddot r}{r} - \frac{2\|q\|_2^2}{r^3}(r^2\dot a) = \ddot b - u\ddot r - \frac{2\|q\|_2^2}{r^3}(r^2\dot a).
\end{align}
since the geodesic equations imply $\ddot r = \|q\|_2^2/r^3.$ Now observe that $\ddot b = \ddot u r + 2\dot u \dot r + u\ddot r$, so 
\begin{align}
    \ddot b - u\ddot r - \frac{2\|q\|_2^2}{r^3}(r^2\dot a) &= \ddot u r  + 2\dot u \dot r - \frac{2\|q\|_2^2}{r^3}(r^2\dot a) \\
    &= \frac{1}{r}(r^2\dot u)' - \frac{2\|q\|_2^2}{r^3}(r^2\dot a).
\end{align}
Thus,
\begin{equation}
    (r^2\dot u)' - \frac{2\|q\|_2^2}{r^2}(r^2\dot a) = 0.
\end{equation}
We obtain the lemma statement by taking $p = r^2\dot a$.  
\qed{}

\begin{restatable}[Time varying oscillator]{lemma}{} \label{lemma: time varying oscillator}
    The ordinary differential equation
    \[\ddot u_{\perp} + \frac{\|q\|_2^2}{r^4}u_\perp = 0\]
    with initial conditions $u_\perp(0)$ and $\dot u_\perp(0)$ satisfies 
    \[u_{\perp}(t) = r(t)\left[\frac{u_{\perp,0}}{r_0}\cos(\|q\|_2\sqrt{2}\tau(t)) + \frac{\dot u_{\perp,0} r_0 - u_{\perp,0} \dot r_0}{\|q\|_2\sqrt{2}}\sin(\|q\|_2\sqrt{2}\tau(t))\right].\]
\end{restatable}
\noindent \textit{Proof.} We want to solve the time-varying harmonic oscillator,
\begin{equation}
    \ddot x + \omega^2(t)x = 0
\end{equation}
where $\omega(t) = \|q\|_2/r^2(t).$ We will start by introducing a complex \textit{ansatz}, and then massaging it to solve our desired ODE. Let
\begin{align}
    \chi(t) = \rho(t)e^{i\theta(t)}
\end{align}
where $\rho >0$ is a real-valued amplitude and $\theta$ is a real-valued phase. We have
\begin{equation}
    \dot \chi = \dot \rho e^{i\theta} + i\dot\theta\rho e^{i\theta}
\end{equation}
and 
\begin{equation}
    \ddot \chi = \left(\ddot \rho  + 2i\dot \theta\dot\rho + i\ddot \theta\rho -\dot\theta^2\rho  \right)e^{i\theta}.
\end{equation}
Plugging into the oscillator and canceling $e^{i\theta}$ yields,
\begin{align}
    (\ddot \rho -\dot\theta^2\rho + \omega^2(t)
    \rho)  + i(2\dot \theta\dot\rho + \ddot \theta\rho) = 0.
\end{align}
This equation is satisfied when both the real and complex parts are null. For the complex component,
\begin{align*}
    2\dot \theta\dot\rho + \ddot \theta\rho &= 0 \\
    \iff (\rho^2\dot\theta)' &= 0 \\
    \iff \rho^2\dot\theta &= C 
\end{align*}
for some $C$ constant in time. So we'll choose $C = \sqrt{2}$, implying $\dot\theta = \sqrt{2}/\rho^2$. Plugging into the real component,
\begin{align*}
    \ddot \rho -\frac{2}{\rho^4}\rho + \omega^2(t)
    \rho &= 0 \\
    \iff\ddot \rho  + \omega^2(t)
    \rho &= \frac{2}{\rho^3}.
\end{align*}
This is an example of an Ermakov-Pinney equation \cite{ermakov1880equations, pinney1950nonlinear}. It turns out that choosing $\rho  = \alpha r$ for some $\alpha > 0$ solves this equation. Verifying,
\begin{align}
    \alpha\ddot r + \frac{\|q\|_2^2}{r^3}\alpha = \frac{2}{\alpha^3r^3}.
\end{align}
The geodesic equations of \cref{eq: geodesic equations on the cone} indicate that $\ddot r = r\|v\|^2_2 = \|q\|_2^2/r^3$ since we have defined $q = r^2v$. Thus,
\begin{align}
    2\frac{\alpha\|q\|_2^2}{r^3} &= \frac{C^2}{\alpha^3r^3} \\
    \implies \alpha &= \frac{1}{\sqrt{\|q\|_2}}
\end{align}
solves the equation. Thus,
\[\rho = \frac{r}{\sqrt{\|q\|_2}}\]
which means
\[\theta(t) = \|q\|_2\sqrt{2}\int_0^t\frac{ds}{r(s)^2} = \|q\|_2\sqrt{2}\tau(t)\]
where $\theta(0) = 0$ by choice. Now the general solution to the original oscillator can be written as 
\begin{align*}
    x(t) &= A\rho(t)\cos\theta(t) + B\rho(t)\sin\theta(t) \\
    &= A\frac{r(t)}{\sqrt{\|q\|_2}}\cos(\|q\|_2\sqrt{2}\tau(t)) + B\frac{r(t)}{\sqrt{\|q\|_2}}\sin(\|q\|_2\sqrt{2}\tau(t)).
\end{align*}
One can solve for $A,B$ with initial conditions $x(0), \dot{x}(0)$ to obtain
\begin{align*}
    A &= \frac{x_0}{r_0}\sqrt{\|q\|_2} \\
    B &= \frac{\dot x_0 r_0 - x_0\dot r_0}{\sqrt{2\|q\|_2}}
\end{align*}
yielding
\begin{equation}
    x(t) = r(t)\left[\frac{x_0}{r_0}\cos(\|q\|_2\sqrt{2}\tau(t)) + \frac{\dot x_0 r_0 - x_0 \dot r_0}{\|q\|_2\sqrt{2}}\sin(\|q\|_2\sqrt{2}\tau(t))\right].
\end{equation}

\qed{}

{\bibliography{main.bib}}

\end{document}
